% Generated by build_report.py; shared task-level coverage table.
\section{Related works: mức bao phủ S1–S10}
Khảo sát giữ cửa sổ \textbf{21/09/2023–21/09/2026}. Bảng đối chiếu từng paper với mười nhiệm vụ của ta; mức bao phủ phải đọc cùng giả định và giới hạn ở cột cuối.

S1: vị trí hiện tại; S2: nơi dừng; S3: đoạn đường đã đi; S4: liên kết danh tính giữa phiên; S5: cạnh đường kế tiếp; S6: đích chưa tới; S7: nội dung/ý định truy vấn; S8: định vị nhờ dữ liệu người đi cùng; S9: điểm đầu bị che; S10: điểm cuối bị che của chuyến đã hoàn tất.

\textbf{Fully (F):} paper trực tiếp xử lý mục tiêu suy luận của scenario, trong giả định của paper. \textbf{Partially (P):} chỉ xử lý một phần hoặc bối cảnh hẹp hơn. \textbf{CX:} chưa đủ nguồn để phân loại. \textbf{—:} chưa tìm thấy bằng chứng xử lý trực tiếp; không kết luận cơ chế chắc chắn thất bại. F/P là đối chiếu của ta, không phải nhãn do tác giả công bố, không có nghĩa bảo vệ 100\% hoặc đã vượt các ca A/B/C của ta.

\begingroup\small
\begin{longtable}{@{}P{3.387cm}P{1.742cm}P{2.613cm}P{8.614cm}@{}}
\toprule
\textbf{Paper} & \textbf{Fully} & \textbf{Partially / CX} & \textbf{Căn cứ và giới hạn}\\\midrule
\endfirsthead
\toprule
\textbf{Paper} & \textbf{Fully} & \textbf{Partially / CX} & \textbf{Căn cứ và giới hạn}\\\midrule
\endhead
\bottomrule\endfoot
TransProtect 2024 [\href{https://arxiv.org/html/2409.09495v1}{R1}] & S1, S3 & S2 & §3–5: suy vị trí từ chuỗi đã làm nhiễu, biết đường/traffic; chưa đánh giá riêng nơi dừng.\\
\addlinespace[3pt]
Semantic correlation 2026 [\href{https://link.springer.com/article/10.1007/s44443-026-00899-w}{R2}] & S1, S3 & S2 & §3–6: giữ dummy hợp lý theo thời gian/ngữ nghĩa; chưa chấm riêng nơi dừng hay ý định S7.\\
\addlinespace[3pt]
Fake queries 2026 [\href{https://link.springer.com/article/10.1007/s44443-025-00438-z}{R3}] & S1, S3 & S2 & §3–6: chèn truy vấn giả để khó nối tuyến; chưa đánh giá riêng tọa độ nơi dừng.\\
\addlinespace[3pt]
Road-aware PTPPM 2025 [\href{https://arxiv.org/html/2511.21020v1}{R4}] & S1, S3 & S2 & §III, VI: chống suy vị trí/quỹ đạo có tương quan; bước dự báo chưa là phép thử S5/S6.\\
\addlinespace[3pt]
CPCROK 2025 [\href{https://uhra.herts.ac.uk/id/eprint/25705/1/futureinternet-17-00165.pdf}{R5}] & — & S4 & §3–5: chống nối bí danh xe trong VANET; chưa tách danh tính người/thiết bị như S4.A/B/C.\\
\addlinespace[3pt]
DP-FETC 2025 [\href{https://www.aimspress.com/aimspress-data/era/2025/11/PDF/era-33-11-293.pdf}{R6}] & — & S3, S8, S9, S10 & SynFE/PreOD: sinh dữ liệu offline, đổi đầu/cuối của tuyến tương quan; không chấm định vị từ người đi cùng.\\
\addlinespace[3pt]
Improved PIR 2025 [\href{https://onlinelibrary.wiley.com/doi/abs/10.1002/cpe.70447}{R7}] & CX & CX: S1, S7 & Tóm tắt nêu bảo vệ vị trí/nội dung bằng mã hóa; thiếu toàn văn để đối chiếu quyền quan sát và suy ý định.\\
\addlinespace[3pt]
SoK 2024 [\href{https://petsymposium.org/popets/2024/popets-2024-0068.pdf}{R8}] & Không áp dụng & Không áp dụng & Khảo sát cách đánh giá/tấn công; không phải cơ chế bảo vệ để gán F/P.\\
\addlinespace[3pt]
PRISM 2025 [\href{https://www.sciencedirect.com/science/article/pii/S1084804525001687}{R9}] & CX & CX: S1, S3 & Tóm tắt: bảo vệ chia sẻ quỹ đạo; chưa đủ mô hình attacker. Dùng LSTM không xác nhận xử lý S5/S6.\\
\addlinespace[3pt]
Options to Action 2026 [\href{https://elathan.github.io/papers/chi26.pdf}{R10}] & Không áp dụng & Không áp dụng & Khảo sát sử dụng tính năng, có che đầu/cuối (S9/S10); không đo khả năng chống suy luận.\\
\addlinespace[3pt]
Forsch et al. 2023 [\href{https://link.springer.com/chapter/10.1007/978-3-031-35374-1_3}{R11}] & S9, S10 & — & §3.2: S-TT cắt cả hai đầu, xét nhiều tuyến cùng đích; bảo đảm theo tập tấn công/địa điểm cho trước, offline.\\
\addlinespace[3pt]
EV querying 2024 [\href{https://wrap.warwick.ac.uk/id/eprint/183198/1/WRAP-privacy-preserving-querying-mechanism-high-utility-electric-vehicles-2024.pdf}{R12}] & S1, S3 & S2 & Proposed Query Model: AGeoI + dummy chống định vị/nối tuyến online; chưa chấm riêng nơi dừng hoặc S4–S10.\\
\end{longtable}
\endgroup
Scenario không được liệt kê ở một hàng được hiểu là —; với R7/R9 là chưa xác minh, với R8/R10 là không áp dụng. Bảng cho thấy các cơ chế chuỗi chủ yếu xét S1/S3; S2 cần phép thử nơi dừng riêng, còn S9/S10 cần kiểm tra suy ngược từ phần tuyến vẫn công bố.

